\minitopic{模型输出是什么认识对象}生成模型直接产出的是符号序列，不是带有已验证来源的事实数据库记录。序列可准确复述、合理推断、创造性改写，也可在相同流畅度下虚构。用户不能从语气判断属于哪一种。认识评价应落在具体人机流程，而非抽象问“模型总体可靠吗”。 同一模型在改写用户提供文本时，可通过逐句对照验证；在回答冷门历史细节时，训练来源与事实对应不可见；在数学推导时，可用形式检查；在价值建议时，问题又涉及目标。任务类型决定验证方法和可接受依赖。 模型还可能通过工具检索或计算。工具调用扩大能力，却使信息链更复杂：检索是否找到权威页面，摘要是否忠实，计算参数是否正确，最终语言是否保持结果。显示“已联网”不是认识保证，必须保留每步可追溯证据。

\minitopic{虚构与传统错误的差异}NIST把生成式AI虚构定义为自信陈述错误或虚假内容，并把它与有害偏差、人机配置、信息完整性等风险区分\parencite{nist2024}。虚构不是模型拥有错误记忆后有意欺骗，也不等同于普通打字错误；它源于生成目标并不内含逐项事实验证。 “幻觉”一词容易拟人化，仿佛系统具有知觉经验。本书优先用“虚构”描述输出性质，同时不把所有错误都归为虚构。用户提示含错前提、外部数据库过时、检索解析失败和计算工具故障具有不同成因与对策。 错误分类影响责任。若模型在无来源任务中虚构，部署者应提供核验界面和限制；若权威数据库自身错误，需修正数据治理；若用户忽略明确警告，责任部分转向使用；若系统用肯定语气掩盖不确定，界面设计不能免责。

\minitopic{流畅性与认识错觉}人会把易处理、熟悉和重复的信息判断得更真。生成文本语法完整、结构均衡、即时响应，提升处理流畅性，也可能制造“已经理解”的感觉。流畅不是坏事，教材本就应清楚；风险在于形式质量与证据质量脱钩。 引用样式也能制造权威感。一个含作者、年份和DOI的虚构条目比含糊陈述更像知识。核验必须把题名、作者、期刊、年份、卷期、页码和标识符分别检查，而不是链接能打开就算通过。链接可能真实却不支持所引主张。 用户还会因系统能解释而高估自己理解。真正理解应能脱离原答案重构、回答变体、发现限制。学习者可在看输出前先作答，随后比较；或让系统提出反例，再由原典裁决。工具应增加认知活动而非替代活动。

\minitopic{检索增强的认识边界}检索增强生成把外部文档加入上下文，能降低部分无来源虚构，却不能保证忠实。检索排序可能漏掉反证，文档可能过时，生成器可能把两个来源混合或作超出文本的推断。需要同时评价召回、来源质量和蕴涵忠实度。 引用覆盖率是一个有用指标：答案中哪些可核验主张有来源？引用正确性则问来源是否真的支持。即使两者都高，答案可能遗漏关键观点或把价值判断包装成事实。完整性与相关性仍需人工或独立程序评价。 在高风险领域，检索库应版本化，记录访问时间与文档身份。否则同一问题后来返回不同来源，无法重现决策。可重复性是人机知识链的制度条件。

\minitopic{训练数据与来源问题}模型从大量材料学习统计结构，通常不能为某个输出逐句指出训练来源。不能把“训练中可能见过”当作引证。知识产权、隐私和事实溯源也是不同问题：一段输出不侵犯版权，不等于事实可靠；事实正确，也不等于数据使用正当。 训练数据代表性影响性能分布。主流语言和高资源领域可能更稳定，低资源语言、地方知识与边缘群体记录更稀缺。平均基准会遮蔽这种差异。部署评估应按实际用户、语言和任务分层。 数据规模也会复制错误和刻板印象。更多数据不是自动更客观；来源权重、重复内容和缺失群体决定学到的分布。社会认识论由此进入模型评价：谁的陈述被收集，谁有修正权，错误影响由谁承担？

\minitopic{内容溯源的能力与限度}C2PA内容凭证用加密签名与清单记录数字资产来源和修改历史，帮助验证溯源信息是否被篡改\parencite{c2pa2025}。它回答“谁以何工具在何链条中声明了什么”，而不单独回答图像所描绘事件是否真实。 签名真实不保证签名者诚实，完整编辑记录不保证内容语境正确；没有凭证也不证明内容虚假，因为工具可能不支持或元数据被移除。溯源是证据之一，必须与来源信誉、现场信息和内容分析结合。 凭证还带来隐私与权力问题。强制身份可暴露记者、举报者和弱势创作者；可选择披露保护隐私，却降低某些核验。标准设计需在可追溯、匿名和滥用风险之间权衡。

\minitopic{自动化偏差的机制}人可能因机器建议而停止搜索、把建议设为默认，或在冲突时不当地改变正确判断。自动化偏差不只是“太信机器”，还包括工作流把拒绝设计得困难、时间压力使复核成为形式。把人放在环中只有在其有信息、时间、权限和替代方案时才有效。 算法厌恶则在看到一次机器错误后过度放弃，即使人类基准更差。两种偏差可并存：人对熟悉系统过信，对陌生系统过度拒绝。校准应基于分任务错误记录，而非品牌印象或单个轶事。 分歧界面可以提升认识：明确显示人机不同、要求先独立判断、突出模型适用边界。若系统总先给答案，会锚定人；若只在低信心时请人处理，又可能把最难案例全部交给缺少背景的人。

\minitopic{解释的多重目标}“模型可解释”可能指工程师能调试、受影响者能质疑、医生能判断临床合理、监管者能审计，或普通用户能形成恰当信任。一个全局特征重要性图不能同时满足所有目标。 解释也可能忠实或仅可理解。事后生成一个简洁理由若不反映实际决策机制，会成为新的误导文本。复杂真实机制难以沟通，简化又可能失真。应说明解释对象、受众与保真程度。 认识上最重要的解释有时是能力说明：系统在哪类数据评估过，错误如何分布，遇到分布变化会怎样，谁监控。因果内部解释并非信任的唯一基础，外部验证也不可缺少。

\minitopic{AI与认识劳动的重新分配}自动化可减少搜索和格式劳动，让人投入问题设计与批评；也可使初级研究者失去学习机会，把能力集中在少数系统提供者。短期效率与长期能力形成可能冲突。教育使用应问：哪些步骤本身就是学习目标？ 若学生用系统生成论证后只做表面润色，最终作品难证明其能力。若先独立重构，再用系统生成反例并逐项核验，工具可扩大练习。评价应关注过程记录、口头迁移和来源审计，而非只做AI文本侦测。 机构依赖专有系统还会形成认识脆弱性：模型更新、服务中断和不可审计改变知识链。关键流程需要版本记录、替代渠道与退出方案。

\minitopic{机器知识的四种立场}行为立场以系统能正确回答、推理和行动为准，较容易承认机器知识。功能立场要求内部状态在系统中发挥信念与证据作用。外在主义立场可让可靠表征在适当环境构成知识。主体性立场则要求理解、理由响应、承诺或责任，对现有系统更谨慎。 争论应避免从词语直接推本体。系统输出“我知道”不证明有知识，系统没有人类意识也不直接证明任何功能知识不可能。先说明知识理论，再把系统架构和能力代入。 实践中，人机组合问题更紧迫。即使模型自身不知，医生可否借其获得知识？如果模型通过独立验证、错误可监控、医生理解适用边界，答案可能肯定；若输出不可追溯且医生盲从，组合不可靠。

\minitopic{综合案例：生成式法律检索}律师让系统寻找支持主张的判例。系统给出三条，其中两条真实但已被后案限制，一条完全虚构。律师若直接提交，流畅性与格式掩盖了来源失败。即使法律结论碰巧正确，也缺少适当根据并可能违反专业责任。 合格流程包括在官方数据库核对案件、审查后续效力、阅读关键段落、保存检索日期，并让另一人复核高风险引用。系统可用于生成搜索词、整理已验证材料，而不是作为最终权威。 若未来系统直接连接权威数据库并给逐段链接，风险降低但不消失：检索可能漏掉不利判例，摘要可能越界，法律适用含价值与类比判断。技术改进改变核验成本，不取消人的论证责任。

\begin{tcolorbox}[breakable,colback=white,colframe=blue!30!black,title={人机认识链审计表},fonttitle=\bfseries]
标明任务类型、模型版本与工具；把输出拆成事实、推断和建议；逐项核对权威来源及蕴涵；检查分层性能与错误机制；记录人是否有时间、权限和能力复核；区分内容溯源与内容真实性；保存版本、提示、来源和最终人工决定。
\end{tcolorbox}
